\documentclass[10pt,a4paper]{amsart}
\usepackage{amsmath,amssymb,amsthm,graphicx}
\usepackage{verbatim}
\newenvironment{reference}{\comment}{\endcomment}
%\newenvironment{reference}{}{}
\newenvironment{slogan}{\comment}{\endcomment}
\newenvironment{history}{\comment}{\endcomment}
% For commutative diagrams we use Xy-pic
\usepackage[all]{xy}

% We use 2cell for 2-commutative diagrams.
\xyoption{2cell}
% Theorem environments.
%
\theoremstyle{plain}
\newtheorem{theorem}[subsection]{Theorem}
\newtheorem{proposition}[subsection]{Proposition}
\newtheorem{lemma}[subsection]{Lemma}

\theoremstyle{definition}
\newtheorem{definition}[subsection]{Definition}
\newtheorem{example}[subsection]{Example}
\newtheorem{exercise}[subsection]{Exercise}
\newtheorem{situation}[subsection]{Situation}

\theoremstyle{remark}
\newtheorem{remark}[subsection]{Remark}
\newtheorem{remarks}[subsection]{Remarks}

\numberwithin{equation}{subsection}

% Macros
%
\def\lim{\mathop{\mathrm{lim}}\nolimits}
\def\colim{\mathop{\mathrm{colim}}\nolimits}
\def\Spec{\mathop{\mathrm{Spec}}}
\def\Hom{\mathop{\mathrm{Hom}}\nolimits}
\def\Ext{\mathop{\mathrm{Ext}}\nolimits}
\def\SheafHom{\mathop{\mathcal{H}\!\mathit{om}}\nolimits}
\def\SheafExt{\mathop{\mathcal{E}\!\mathit{xt}}\nolimits}
\def\Sch{\mathit{Sch}}
\def\Mor{\mathop{\mathrm{Mor}}\nolimits}
\def\Ob{\mathop{\mathrm{Ob}}\nolimits}
\def\Sh{\mathop{\mathit{Sh}}\nolimits}
\def\NL{\mathop{N\!L}\nolimits}
\def\CH{\mathop{\mathrm{CH}}\nolimits}
\def\proetale{{pro\text{-}\acute{e}tale}}
\def\etale{{\acute{e}tale}}
\def\QCoh{\mathit{QCoh}}
\def\Ker{\mathop{\mathrm{Ker}}}
\def\Im{\mathop{\mathrm{Im}}}
\def\Coker{\mathop{\mathrm{Coker}}}
\def\Coim{\mathop{\mathrm{Coim}}}

% Boxtimes
%
\DeclareMathSymbol{\boxtimes}{\mathbin}{AMSa}{"02}

%
% Macros for moduli stacks/spaces
%
\def\QCohstack{\mathcal{QC}\!\mathit{oh}}
\def\Cohstack{\mathcal{C}\!\mathit{oh}}
\def\Spacesstack{\mathcal{S}\!\mathit{paces}}
\def\Quotfunctor{\mathrm{Quot}}
\def\Hilbfunctor{\mathrm{Hilb}}
\def\Curvesstack{\mathcal{C}\!\mathit{urves}}
\def\Polarizedstack{\mathcal{P}\!\mathit{olarized}}
\def\Complexesstack{\mathcal{C}\!\mathit{omplexes}}
% \Pic is the operator that assigns to X its picard group, usage \Pic(X)
% \Picardstack_{X/B} denotes the Picard stack of X over B
% \Picardfunctor_{X/B} denotes the Picard functor of X over B
\def\Pic{\mathop{\mathrm{Pic}}\nolimits}
\def\Picardstack{\mathcal{P}\!\mathit{ic}}
\def\Picardfunctor{\mathrm{Pic}}
\def\Deformationcategory{\mathcal{D}\!\mathit{ef}}
\usepackage{cleveref}
\crefname{originallemma}{original lemma}{original lemmas}
\Crefname{originallemma}{Original lemma}{Original lemmas}
\makeatletter
\newlabel{lemma-affine-base-change}{{07U8}{source}{}{sourcepointer-1}{}}
\newlabel{lemma-affine-base-change@cref}{{[originallemma][0][]07U8}{[1][0][]source}{}{}{}}
\makeatother
\title{The counit of a quasi-compact separated etale morphism of algebraic spaces is a split surjection}
\author{The Stacks Project Authors}
\date{stacks-project:spaces-cohomology.tex:local-snapshot-20261012; GNU FDL 1.2 or later}
\begin{document}
\maketitle
\noindent\textit{Editorial note.} All delivered mathematical text is the author's own, copied from the pinned Stacks Project chapter file spaces-cohomology.tex; no statement or proof step was rewritten, reordered or removed.\par\medskip
\noindent\textit{Editorial note.} Preamble provenance: the pinned chapter file opens with \textbackslash{}input\{preamble\} and carries no preamble text of its own, so the project's single Stacks-wide preamble staged at sources/stacks-spaces-cohomology/preamble.tex is excerpted verbatim as three spans, lines 8-12 (verbatim package and the project's own reference/slogan/history display environments), lines 14-18 (Xy-pic setup) and lines 154-217 (theorem declarations and mathematical macros). The wrapper class and the amsmath/amssymb/amsthm/graphicx base packages are supplied by the assembler. No macro was invented or redefined.\par\medskip
\noindent\textit{Editorial note.} Same-chapter pointers are located at their \textbackslash{}label lines inside the pinned chapter file; cross-chapter pointers are located in the project's own tag list staged in the snapshot (sources/stacks-chow/tags.txt), at the verbatim lines carrying their canonical tags.\par\medskip
\noindent\textit{Allowed standard prerequisites (editorial).}\begin{itemize}
\item The projection formula identifying the pullback of a pushforward along the fibre square. Same chapter spaces-cohomology.tex, label lemma-affine-base-change (invoked at proof line 2036). Applied to the cartesian square with the two projections \$p\$ and \$q\$, it gives \$f\textasciicircum{}*f\_*\textbackslash{}mathcal\{F\} = q\_*p\textasciicircum{}*\textbackslash{}mathcal\{F\}\$, which is the sheaf whose relation to \$\textbackslash{}mathcal\{F\}\$ is being described.
\item The diagonal of a separated etale morphism is an open and closed immersion. Standard: \$\textbackslash{}Delta : X \textbackslash{}to X \textbackslash{}times\_Y X\$ is an open immersion because \$f\$ is etale and a closed immersion because \$f\$ is separated (used at proof lines 2038-2039). It splits the structure sheaf of the fibre square into the diagonal and the complementary open and closed piece, which is what makes \$p\textasciicircum{}*\textbackslash{}mathcal\{F\}\$ decompose.
\end{itemize}\medskip
\ifdefined\hypertarget\hypertarget{sourcepointer-1}{}\fi
\noindent\textit{Source reference.} Original lemma 07U8: sources/stacks-spaces-cohomology/pinned-original-source.tex, source lines 1902--1902. Navigation pointer for \textbackslash{}ref\{lemma-affine-base-change\} used inside the counted span (line 2036); the argument that consumes it is retained in the span.\par\medskip
\begin{lemma}
\label{lemma-etale-pull-push-split}
Let $f : X \to Y$ be a quasi-compact, separated, \'etale morphism
of algebraic spaces. Then for any quasi-coherent $\mathcal{O}_X$-module
$\mathcal{F}$ the map $f^*f_*\mathcal{F} \to \mathcal{F}$ is split.
\end{lemma}

\begin{proof}
Consider the cartesian diagram
$$
\xymatrix{
X \times_Y X \ar[r]_-p \ar[d]_q & X \ar[d]^f \\
X \ar[r]^f & Y
}
$$
By Lemma \ref{lemma-affine-base-change} we have
$f^*f_*\mathcal{F} = q_*p^*\mathcal{F}$. The morphism
$\Delta : X \to X \times_Y X$ is open (as a morphism between
algebraic spaces \'etale over $Y$) and closed as $f$ is separated.
Thus we see that $p^*\mathcal{F}$ is a direct sum of
$\Delta_*\mathcal{F}$ and a quasi-coherent module supported
on the (open and closed) complement of $\Delta(X)$.
Tracing the maps the reader verifies that
$$
\mathcal{F} = q_*\Delta_*\mathcal{F} \to
q_*p^*\mathcal{F} = f^*f_*\mathcal{F} \to
\mathcal{F}
$$
is the identity map. Details omitted.
\end{proof}

\end{document}
